\documentclass[10pt,letterpaper,final]{report}
\usepackage[margin=.75in]{geometry}
\usepackage[utf8]{inputenc}
\usepackage{amsmath}
\usepackage{amsfonts}
\usepackage{amssymb}
\usepackage{amsthm}
\renewcommand\qedsymbol{$\blacksquare$}
\usepackage{enumerate}
\usepackage{enumitem}
\usepackage{hyperref}
\usepackage{pdfpages}
\usepackage{graphics}
\usepackage{graphicx}
\usepackage{tikz}
\usepackage{tikz-qtree}
\usepackage{forest}
\usepackage{float}
\usetikzlibrary{automata,arrows}
\usetikzlibrary{shapes.multipart, positioning}
\usepackage{mdframed}
\usepackage[
  backend=biber,
  style=numeric-comp,
  sorting=none
]{biblatex}
\addbibresource{references.bib}

\usepackage{minted}
\usemintedstyle{algol_nu}
\usepackage{xltabular}
\usepackage{pdflscape}
\usepackage{tcolorbox}
\usepackage{array}
\usepackage{todonotes}
\setuptodonotes{inline}  % render \todo notes as full-width blocks in the text flow, not margin boxes

\newcommand{\reason}[1]{&& \text{#1}}

% Based on template from COSC-417 at Towson University, originally authored by Marius Zimand

% To input a script, adapt the following:
% \inputminted[breaklines, breakanywhere, fontsize=\footnotesize, frame=single, fontsize=\footnotesize]{python}{script.py}

\begin{document}
\begin{center}
  \fbox{
    \vbox{
      \begin{flushleft}
        Jonathan Llovet \\  % authors' names
        EN.625.252.81 \\  %class
        2026-10-11 \\  % date
        Module 6 - Homework 6 \\ % ID
      \end{flushleft}
      \center{\Large{\textbf{Linear Algebra}}}
      %\end{mdframed}
    } % end vbox
  } % end fbox
  \vline
\end{center}

\medskip
\noindent\textbf{Problem 1:}
\medskip

Consider the matrix

\begin{align*}
  \begin{bmatrix}
    1  & 1 & 1  \\
    -1 & 1 & -1 \\
    2  & 0 & 2
  \end{bmatrix}
\end{align*}


\begin{enumerate}[label=(\alph*)]
  \item Compute its rank.
  \item Find a basis for $Col(A)$
  \item Find a basis for $Nul(A)$
  \item State the Rank Theorem from the text.
  \item Show your answers to the previous parts of this question satisfy the Rank Theorem.
\end{enumerate}

\medskip
\noindent\textbf{Answer:}

\pagebreak

\noindent\textbf{Problem 2:}
\medskip

Let $S$ be a subspace of $\mathbb{R}^4$ spanned by the set of vectors

\begin{align*}
  S = \begin{Bmatrix}
        \begin{bmatrix}
          1  \\
          -1 \\
          0  \\
          1  \\
        \end{bmatrix},
        \begin{bmatrix}
          -2 \\
          1  \\
          1  \\
          1  \\
        \end{bmatrix},
        \begin{bmatrix}
          0 \\
          0 \\
          0 \\
          0 \\
        \end{bmatrix},
        \begin{bmatrix}
          -1 \\
          0  \\
          1  \\
          2  \\
        \end{bmatrix},
        \begin{bmatrix}
          1  \\
          1  \\
          -2 \\
          -5 \\
        \end{bmatrix}
      \end{Bmatrix}
\end{align*}

\begin{enumerate}[label=(\alph*)]
  \item Find a basis for Span $S$. Perform all row reduction by hand, showing all your steps. (\emph{Hint:} before you start to row reduce, can you remove any obviously linearly dependent vectors in $S$ to simplify the row reduction?)
  \item State the dimension of Span $S$.
\end{enumerate}

\medskip
\noindent\textbf{Answer:}
\medskip

\pagebreak

\section*{Source Code}
The full source code, including LaTeX, tooling, other artifacts, and git history is available on my Github repository for the class here:
\begin{center}
  \url{https://github.com/jllovet/jhu-en-625-252-fa26-linear-algebra}
\end{center}

\printbibliography

\end{document}
